\begin{enabstract}

In social media, information propagation is deeply influenced by users' cognitive processes. Upon receiving information, users must undergo cognitive processing stages such as attention capture, content comprehension, and attitude formation before engaging in propagation behaviors like replying and reposting. This cognition-driven propagation process is referred to as information cognitive propagation. In recent years, large language model-powered AI agents, represented by Grok, have begun independently participating in discussions on the X platform. Human users and AI agents, as entities with fundamentally different cognitive and content generation mechanisms, now coexist within the same propagation ecosystem. However, existing information propagation analysis and cascade prediction methods generally assume homogeneity among propagation nodes without incorporating agent type into their modeling frameworks. Moreover, publicly available datasets lack agent type annotations, which constrains research on multi-agent propagation behavior. To address these issues, this thesis designs and implements a multi-agent information cognitive propagation evaluation system that employs propagation behavior metrics as proxy variables for cognitive processes, quantitatively assessing the behavioral differences between human users and AI agents across three evaluation dimensions: the message level, the cascade level, and the interaction type level. The main contributions are as follows:

(1) A Group-aware Fusion TabTransformer model is proposed for comment propagation value assessment. Targeting the message-level evaluation dimension, the model introduces agent type as an explicit feature group and employs a gating fusion mechanism to compute sample-level feature group weights, adaptively adjusting the contributions of four heterogeneous feature groups --- exogenous context, text semantics, sentiment orientation, and agent identity --- to achieve comment-level propagation value quantification.

(2) An Interaction-aware Heterogeneous Path Encoder (IHPE) is proposed for propagation continuation prediction. Targeting the cascade-level and interaction-type-level evaluation dimensions, the model takes the complete path from the root node to the target node in the comment tree as input, explicitly models four types of agent interaction patterns along the path through an interaction type encoding layer, and combines type-augmented node encoding with LSTM-based sequential encoding to achieve path-level propagation continuation prediction.

(3) A complete evaluation system is designed and implemented, adopting a four-layer architecture with six functional modules that encompass data import and preprocessing, model inference, multi-agent comparative analysis, and visualization. The system aggregates the outputs of both models by agent type to systematically quantify propagation behavior differences between the two types of agents. Functional and performance tests verify the correctness and usability of the system.

    \enkeywords{information cognitive propagation; multi-agent; propagation evaluation; deep tabular learning; heterogeneous path encoding}
\end{enabstract}